\documentclass[a4paper]{article}

\usepackage{pdp}
\usepackage{booktabs}
\usepackage[utf8]{inputenc}
\usepackage[T1]{fontenc}
\usepackage[french]{babel}
\usepackage{tikz}
\usepackage{tcolorbox}
\usetikzlibrary{trees,shapes,arrows}

\tcbuselibrary{skins, breakable}
\newtcolorbox{besoinbox}{
  colback=white,
  colframe=blue!75!black,
  fonttitle=\bfseries,
  coltitle=black,
  colbacktitle=blue!75!black,
  enhanced,
  attach boxed title to top left={yshift=-3mm, xshift=5mm},
  boxrule=0.5pt,
  breakable,
  arc=4pt,
  left=6pt,
  right=6pt,
  top=6pt,
  bottom=6pt
}

%------------------------------------------------------------------------------
%	REPORT INFORMATION
%------------------------------------------------------------------------------

\title{ProjectName -- Java}
\subtitle{Rapport final}
\author{Étudiant1, Étudiant2, Étudiant3, Étudiant4, Étudiant5}

%------------------------------------------------------------------------------

\begin{document}
	% Définition des styles TikZ pour les nœuds des schémas
	\tikzstyle{decisionMax} = [isosceles triangle, draw, shape border rotate=90, minimum width=1.5cm, inner sep=2pt, font=\small, align=center]
	\tikzstyle{decisionMin} = [isosceles triangle, draw, shape border rotate=270, minimum width=1.5cm, inner sep=2pt, font=\small, align=center]
	\tikzstyle{chance} = [circle, draw, minimum size=1.2cm, font=\small, align=center]
	\tikzstyle{leaf} = [text width=1cm, text centered, font=\scriptsize]
%------------------------------------------------------------------------------
%	TITLE PAGE
%------------------------------------------------------------------------------

\maketitle
\thispagestyle{empty}

%------------------------------------------------------------------------------
%	PARTIE 1 — ALGORITHMES
%------------------------------------------------------------------------------

	\section{L'Algorithme Minimax :}

	Conçu pour des jeux à information parfaite et à somme nulle comme les Échecs, l'algorithme Minimax modélise l'arbre des coups possibles en alternant un tour où l'IA cherche à maximiser son score, et un tour où elle simule un adversaire cherchant à minimiser l'avantage de l'IA.

	Pour adapter cet algorithme au Scrabble (où le chevalet adverse est inconnu), notre implémentation déduit d'abord les tuiles restantes en soustrayant du jeu complet les lettres déjà posées sur le plateau et celles présentes dans le chevalet de l'IA. Ensuite, le solveur utilise une méthode de Monte-Carlo : il simule l'adversaire en tirant aléatoirement 200 chevalets virtuels de 7 lettres parmi ces tuiles restantes. Le Minimax évalue alors le pire scénario possible : parmi ces 200 simulations, il retient le coup adverse qui marque le plus de points, et le soustrait du score du coup envisagé par l'IA.

	\subsection*{L'impact de l'heuristique}
	Au Scrabble, évaluer un coup uniquement par les points qu'il rapporte sur le plateau n'est pas le plus pertinent. L'algorithme intègre donc une fonction heuristique d'évaluation qui note la qualité des lettres conservées sur le chevalet après un coup. Cette heuristique donne des bonus et malus en fonction de:

	\begin{itemize}
		\item Un joker conservé rapporte un bonus massif de +15.0 points.
		\item La lettre 'S' rapporte +8.0 points.
		\item Les lettres difficiles (J, K, Q, X, Z) appliquent un malus de -6.0 points.
		\item L'équilibre voyelles/consonnes est récompensé (+5.0 points) tandis qu'un chevalet exclusivement composé de voyelles ou de consonnes est lourdement pénalisé (-12.0 points).
	\end{itemize}

	\subsection*{Complexité et coût}
	Dans notre implémentation de Monte-Carlo, l'algorithme ne disposant pas d'élagage Alpha-Beta, la complexité temporelle est strictement identique à celle de l'Expectiminimax, soit $O(B \times N \times B)$ pour une anticipation d'un tour de l'adversaire (où $B$ est le facteur de branchement et $N=200$ le nombre de tirages simulés). La complexité reste modérée car le parcours se fait en profondeur d'abord. Le coût est massif puisqu'il doit évaluer exactement les 200 scénarios complets. L'exécution est donc bridée par une limite de temps stricte (\texttt{timeLimitMs}), l'empêchant de calculer sur plusieurs tours d'anticipation.

	\subsection*{Performances}
	La performance défensive est excellente. En s'attendant systématiquement au pire, l'IA verrouille le plateau et évite d'ouvrir des cases multiplicatrices dangereuses. Mais ce pessimisme absolu la rend parfois trop peu entreprenante sur le plan offensif. De plus, puisqu'elle paie le même coût que l'Expectiminimax sans bénéficier d'une prise de décision nuancée, cette implémentation se révèle, à temps de calcul égal, moins optimale que l'approche probabiliste.

	\section{L'Algorithme Expectiminimax :}

	L'Expectiminimax est l'évolution du Minimax pour les jeux intégrant du hasard. Il introduit un nouveau type de nœud dans l'arbre de recherche : les "nœuds de chance" (chance nodes). Au lieu de considérer que l'adversaire jouera le pire coup absolu, l'algorithme calcule l'espérance mathématique de tous les coups possibles en fonction de leur probabilité.

	Dans notre code, la méthode \texttt{expectiminimax()} utilise la même base de simulation que le Minimax (les 200 tirages aléatoires de chevalets adverses). Cependant, au lieu d'extraire la valeur maximale de ces simulations, l'algorithme calcule la moyenne des meilleurs scores adverses obtenus sur ces 200 échantillons. Il évalue ainsi un coup en soustrayant cette "réponse moyenne attendue" du score du coup de l'IA.

	\subsection*{L'impact de l'heuristique}
L'Expectiminimax ne se contente pas de calculer mathématiquement les risques sur le plateau. Pour juger de la force réelle d'un coup, l'IA additionne deux éléments concrets :
\begin{itemize}
    \item \textbf{Le risque sur le plateau (L'algorithme) :} L'IA vérifie qu'elle ne laisse pas d'opportunités trop faciles et dangereuses à l'adversaire (comme ouvrir l'accès à une case "Mot Compte Triple").
    \item \textbf{La qualité de sa main (L'heuristique) :} L'IA évalue les lettres qu'elle garde sur son chevalet pour son prochain tour.
\end{itemize}
C'est le compromis entre ces deux visions qui rend l'IA véritablement intelligente. Par exemple, si un coup sécurise parfaitement le plateau mais oblige l'IA à conserver une main catastrophique, l'heuristique va agir comme un garde-fou et lourdement pénaliser ce choix. L'IA cherche donc toujours l'équilibre parfait entre "limiter la casse sur le plateau" et "garder de bonnes lettres pour l'avenir".

	\subsection*{Complexité et coût}
	La complexité est de $O(B \times N \times B)$, ce qui provoque une explosion combinatoire. Bien que son coût soit strictement identique à celui de notre Minimax actuel, l'Expectiminimax a l'inconvénient théorique de ne presque pas pouvoir être optimisé par des techniques d'élagage (comme l'Alpha-Beta). En effet, le calcul d'une espérance mathématique exige d'explorer toutes les probabilités jusqu'au bout, ce qui le coupe quasi systématiquement avant d'atteindre une profondeur de recherche de 2.

	\subsection*{Performances}
	Sur le papier et dans notre code actuel, c'est la meilleure performance stratégique pour le Scrabble. L'évaluation du risque est mieux gérée : l'IA n'hésitera pas à ouvrir un "Mot Compte Triple" si elle calcule que la probabilité que l'adversaire possède les lettres exactes pour s'y placer est très faible. Le style de jeu devient plus naturel. Étant donné que son coût est identique à celui du Minimax dans notre code, l'Expectiminimax s'avère supérieur dans la prise de ses décisions.

	% ------------------------------------------------------------------------------
	% Section : Schémas Comparatifs
	% ------------------------------------------------------------------------------
	\section{Comparaison des deux Algorithmes}

	Les schémas ci-dessous mettent en évidence que la différence entre nos deux algorithmes n'est pas structurelle (l'arbre exploré est identique), mais réside uniquement dans la méthode d'agrégation des nœuds de simulation.

	  \begin{figure*}[h!]
		\centering

		% ---------- ARBRE MINIMAX (PESSIMISTE) ----------
		\begin{minipage}{0.48\textwidth}
			\centering
			\begin{tikzpicture}[
				>=stealth,
				edge from parent/.style={draw, ->},
				level 1/.style={sibling distance=3.5cm, level distance=2cm},
				level 2/.style={sibling distance=1.5cm, level distance=2cm},
				level 3/.style={sibling distance=1.5cm, level distance=1.2cm},
				decisionMax/.style={isosceles triangle, draw, shape border rotate=90, minimum height=0.8cm, inner sep=2pt, font=\scriptsize\sffamily},
				decisionMin/.style={isosceles triangle, draw, shape border rotate=270, minimum height=0.8cm, inner sep=2pt, font=\scriptsize\sffamily},
				pessimist/.style={rectangle, draw, minimum size=0.9cm, font=\scriptsize\sffamily},
				leaf/.style={font=\scriptsize\sffamily}
				]

				\node[decisionMax, align=center] {IA\\(MAX)}
				child { node[pessimist, align=center] (min1) {MIN}
					child { node[decisionMin, align=center] {Adv.}
						child { node[leaf] {$s_1$} }
						edge from parent node[left, font=\tiny\sffamily] {Sim 1}
					}
					child { node[font=\scriptsize\sffamily] {\dots} edge from parent[draw=none] }
					child { node[decisionMin, align=center] {Adv.}
						child { node[leaf] {$s_{200}$} }
						edge from parent node[right, font=\tiny\sffamily] {Sim 200}
					}
					edge from parent node[left, font=\scriptsize\sffamily, xshift=-0.1cm] {Coup $m_1$}
				}
				child { node[font=\scriptsize\sffamily] {\dots} edge from parent[draw=none] }
				child { node[pessimist, align=center] {MIN}
					edge from parent node[right, font=\scriptsize\sffamily, xshift=0.1cm] {Coup $m_B$}
				};

				\node[draw, dashed, font=\scriptsize\sffamily, fill=gray!10, align=center] (agg1) at (0, -2.5) {Pire cas\\$\max(s_1..s_{200})$};
				\draw[->, dotted] (agg1) -- (min1);
			\end{tikzpicture}
			\caption{ \textbf{Agrégation Pessimiste}}
		\end{minipage}
		\hfill
		% ---------- ARBRE EXPECTIMINIMAX (PROBABILISTE) ----------
		\begin{minipage}{0.48\textwidth}
			\centering
			\begin{tikzpicture}[
				>=stealth,
				edge from parent/.style={draw, ->},
				level 1/.style={sibling distance=3.5cm, level distance=2cm},
				level 2/.style={sibling distance=1.5cm, level distance=2cm},
				level 3/.style={sibling distance=1.5cm, level distance=1.2cm},
				decisionMax/.style={isosceles triangle, draw, shape border rotate=90, minimum height=0.8cm, inner sep=2pt, font=\scriptsize\sffamily},
				decisionMin/.style={isosceles triangle, draw, shape border rotate=270, minimum height=0.8cm, inner sep=2pt, font=\scriptsize\sffamily},
				chance/.style={circle, draw, minimum size=0.9cm, font=\scriptsize\sffamily},
				leaf/.style={font=\scriptsize\sffamily}
				]

				\node[decisionMax, align=center] {IA\\(MAX)}
				child { node[chance, align=center] (chance1) {Hasard}
					child { node[decisionMin, align=center] {Adv.}
						child { node[leaf] {$s_1$} }
						edge from parent node[left, font=\tiny\sffamily] {Sim 1}
					}
					child { node[font=\scriptsize\sffamily] {\dots} edge from parent[draw=none] }
					child { node[decisionMin, align=center] {Adv.}
						child { node[leaf] {$s_{200}$} }
						edge from parent node[right, font=\tiny\sffamily] {Sim 200}
					}
					edge from parent node[left, font=\scriptsize\sffamily, xshift=-0.1cm] {Coup $m_1$}
				}
				child { node[font=\scriptsize\sffamily] {\dots} edge from parent[draw=none] }
				child { node[chance, align=center] {Hasard}
					edge from parent node[right, font=\scriptsize\sffamily, xshift=0.1cm] {Coup $m_B$}
				};

				\node[draw, dashed, font=\scriptsize\sffamily, fill=gray!10, align=center] (agg2) at (0, -2.5) {Moyenne\\$\frac{1}{200}\sum s_i$};
				\draw[->, dotted] (agg2) -- (chance1);
			\end{tikzpicture}
			\caption{\textbf{Agrégation Probabiliste}}
		\end{minipage}

	\end{figure*}

	\subsection*{Analyse des performances}

    Afin de confronter notre étude théorique à la réalité matérielle, nous avons instrumenté le code de notre solveur (\texttt{MinimaxSolver.java}) pour extraire des métriques physiques en conditions réelles de jeu. Le chronomètre d'interruption (\texttt{timeLimitMs}) était fixé à 5000 millisecondes.

    Le tableau ci-dessous synthétise des échantillons représentatifs de la prise de décision de l'IA au cours d'une partie :

    \begin{table}[h!]
        \centering
        \renewcommand{\arraystretch}{1.2}
        \begin{tabular}{@{}lccccc@{}}
            \toprule
            \textbf{Algorithme} & \textbf{Temps} & \textbf{Statut} & \textbf{Coups IA Évalués} & \textbf{Simulations Adv.} & \textbf{Moy. Sim/Coup} \\
            \midrule
            Expectiminimax & 5000 ms & Timeout & 48 sur 84 (57\%) & 9 497 & 197 \\
            Expectiminimax & 3575 ms & Complet & 33 sur 33 (100\%) & 6 600 & 200 \\
            Minimax MC     & 5001 ms & Timeout & 24 sur 58 (41\%) & 4 680 & 195 \\
            Minimax MC     & 5001 ms & Timeout & 17 sur 41 (41\%) & 3 294 & 193 \\
            Minimax MC     & 2801 ms & Complet & 10 sur 10 (100\%) & 2 000 & 200 \\
            \bottomrule
        \end{tabular}
        \caption{Mesures physiques des algorithmes (Profondeur 2, SAMPLES\_COUNT = 200)}
        \label{tab:metrics}
    \end{table}

    \subsection*{Interprétation des résultats}

    L'analyse de ces logs d'exécution met en évidence trois points cruciaux concernant les limites de notre implémentation :

    \begin{itemize}
        \item \textbf{La limite du facteur de branchement :} Les données prouvent que la capacité de l'IA à évaluer l'arbre complet dépend entièrement de la complexité du plateau à l'instant T. Lorsque le nombre de mots possibles est faible (ex: 10 ou 33 coups), le moteur termine son parcours de Profondeur 2 confortablement avant le temps limite (2.8s et 3.5s). À l'inverse, dès que les choix s'élargissent (41, 58 ou 84 coups), l'IA est systématiquement bridée.
        \item \textbf{Les conséquences du coût élevé :} Dans le pire scénario mesuré (17 coups évalués sur 41), le \textit{Timeout} force l'IA à jouer en ignorant totalement 59\% des opportunités légales présentes sur le plateau, l'empêchant potentiellement de trouver le mot optimal.
        \item \textbf{L'efficacité de la limite de temps :} Le moteur est capable de simuler un volume massif de plateaux adverses (jusqu'à 9 497 en 5 secondes). Lorsque le \textit{Timeout} est déclenché, on observe que la moyenne des simulations tombe légèrement sous la barre des 200 (192, 195, 197). Cela démontre que notre algorithme d'interruption fonctionne avec précision : il avorte la boucle interne de Monte-Carlo au milieu de l'évaluation du $N$-ième coup pour garantir la fluidité du jeu en temps réel.
    \end{itemize}

    Ces mesures empiriques confirment de manière irréfutable notre conclusion : bien que la Profondeur 2 soit mathématiquement atteinte, l'absence d'élagage dans nos approches de Monte-Carlo bride physiquement la rentabilité de l'arbre dès que le facteur de branchement dépasse la trentaine de mots candidats.

	\section{Le Machine Learning :}

	Contrairement aux algorithmes de recherche arborescente, le Machine Learning propose une approche statistique. Notre implémentation utilise un réseau de neurones conçu avec TensorFlow. L'architecture du modèle traite la recherche de mots comme un problème de classification multiclasse:

	\begin{itemize}
		\item \textbf{L'entrée :} Le chevalet du joueur est converti en un tenseur de 26 dimensions, représentant la fréquence de chaque lettre (de A à Z).
		\item \textbf{Les couches cachées :} Deux couches denses de 256 et 512 neurones analysent les motifs et anagrammes.
		\item \textbf{La sortie :} Une couche Softmax fournit une distribution de probabilités sur l'ensemble du dictionnaire (une probabilité par mot existant).
	\end{itemize}

	En jeu, l'agent Java interroge ce modèle pré-entraîné, filtre les mots ayant une probabilité de faisabilité supérieure à 0,1\%, et trie les 100 meilleures prédictions.

	\subsection*{Complexité et coût}
	La complexité est divisée en deux parties. La phase d'entraînement a une complexité linéaire dépendante de la taille du dictionnaire, du nombre d'époques et de la taille des lots. En partie, la complexité d'inférence est constante, soit $O(1)$. Le coût en jeu est faible (quelques millisecondes pour une matrice tensorielle). En revanche, le coût en mémoire vive est élevé en raison du chargement natif des bibliothèques TensorFlow et du graphe du modèle dans l'environnement Java.

	\subsection*{Performances}
	La performance pure est très efficace pour trouver des anagrammes complexes. Mais la performance de jeu est nulle de manière autonome : le réseau actuel ignore totalement l'état du plateau et les contraintes de placement. Le moteur doit d'ailleurs prévoir une redirection vers les algorithmes arborescents si aucun des mots prédits par l'IA ne peut être légalement posé.

	\section*{Conclusion :}

	La comparaison de ces méthodes met en évidence qu'aucun algorithme n'est parfait s'il est utilisé de manière isolée. L'Expectiminimax est le champion de la stratégie probabiliste, mais son coût le rend peu intéressant pour des recherches profondes. Le Minimax classique est efficace,surtout avec une bonne heuristique, mais excessivement défensif face à l'inconnu, et notre implémentation le prive de son avantage temporel théorique. Le Machine Learning est un très bon calculateur d'anagrammes, mais il est "aveugle" aux règles de placement sur le plateau.

	La solution la plus efficace réside dans la fusion de ces approches. La complexité de l'Expectiminimax provient de son facteur de branchement immense (générer tous les mots possibles). En utilisant le Machine Learning comme un outil d'élagage probabiliste, le modèle TensorFlow peut instantanément suggérer les 20 meilleurs mots possibles depuis le chevalet. L'Expectiminimax n'a plus qu'à évaluer ces 20 branches réduites face aux risques du plateau, obtenant ainsi la certitude mathématique et probabiliste tout en respectant un temps de calcul extrêmement restreint.

%------------------------------------------------------------------------------
%	PARTIE 2 — SPÉCIFICATIONS : F21 SAUVEGARDE ET RESTAURATION
%------------------------------------------------------------------------------

\part{Spécifications Étendues}

\section{F21 — Sauvegarde et restauration d'une partie}

\subsection*{Description}

Le programme doit pouvoir sérialiser une partie en cours dans un fichier texte ASCII et reconstruire une partie complète à partir d'un fichier de sauvegarde. Le parseur doit être robuste aux erreurs de format et les signaler avec précision (type d'erreur et numéro de ligne). L'implémentation repose sur les classes \texttt{SaveManager.java} et \texttt{GameLoader.java}.

\subsection*{Prérequis}

\begin{itemize}
    \item \textbf{F2 — Configuration :} Paramètres par défaut (langue, blitz, etc.) via \texttt{.scrabblerc} — \texttt{ConfigLoader.java}
    \item \textbf{F10 — Règles du jeu :} État complet du jeu : plateau, scores, joueurs, historique, tour courant
    \item \textbf{F15 — Shell interactif :} Commandes \texttt{save FILE} et \texttt{load FILE} — \texttt{GameControllerAux.java}
    \item \textbf{F22 — Commentaires :} Gestion des commentaires \texttt{\#} et blocs \texttt{\{ \}} — \texttt{GameLoader.java}
    \item \textbf{F23/F24/F25 — Formats :} Sections \texttt{[settings]}, \texttt{[game]} et \texttt{[history]} du fichier de sauvegarde
\end{itemize}

\subsection*{Type de besoin}

Fonctionnel.

% ------------------------------------------------------------------------------
\subsection*{Sous besoins 1. Sauvegarder la partie dans un fichier}
\begin{besoinbox}
\textbf{Description :}
Sérialiser l'état courant du jeu dans un fichier texte ASCII en écrivant successivement les trois sections obligatoires.

\textbf{Fonction associée :}
\begin{verbatim}
void saveGame(Game game, String filePath) throws IOException
\end{verbatim}

\textbf{Actions :}
\begin{itemize}
    \item Ouvrir (ou créer) le fichier au chemin indiqué.
    \item Ecrire successivement \texttt{[settings]}, \texttt{[game]} puis \texttt{[history]}.
    \item Fermer le fichier proprement.
    \item Remonter \texttt{IOException} en cas d'échec (gérée par CLI ou GUI).
\end{itemize}

\textbf{Resultat attendu :}
\begin{itemize}
    \item \textit{Succes :} fichier complet et bien forme ecrit sur le disque.
    \item \textit{Echec :} exception \texttt{IOException} remontee, aucun fichier partiel conserve.
\end{itemize}

\textbf{Test unitaire :}
\begin{itemize}
    \item \textbf{Cas 1 :} Entree : etat complexe (plateau avec mots, 2 joueurs, scores, racks, coups pass/play). Attendu : fichier cree avec 3 sections, verifie via \texttt{SaveManagerTest.java}.
    \item \textbf{Cas 2 :} Entree : conversion de coordonnees \texttt{a1v} / \texttt{h8h}. Attendu : format des coordonnees correct dans \texttt{[history]}.
    \item \textbf{Cas 3 :} Entree : partie vide (aucun coup joue). Attendu : plateau serialise en lignes de tirets, section \texttt{[history]} vide.
\end{itemize}
\end{besoinbox}

\subsection*{Sous besoins 2. Serialiser la section \texttt{[settings]}}
\begin{besoinbox}
\textbf{Description :}
Ecrire les parametres globaux de la partie et les parametres IA joueur par joueur sous la section \texttt{[settings]}.

\textbf{Parametres ecrits :}
\begin{itemize}
    \item \texttt{blitz} — mode blitz active ou non.
    \item \texttt{super-scrabble} — plateau 21$\times$21.
    \item \texttt{players-count} — nombre de joueurs.
    \item \texttt{debug} / \texttt{verbose} — niveaux de trace.
    \item \texttt{player-N-type} — \texttt{human} ou \texttt{ai}.
    \item \texttt{player-N-ai-mode} — algorithme IA (ex.: \texttt{minimax}).
    \item \texttt{player-N-name} — nom du joueur.
\end{itemize}

\textbf{Resultat attendu :}
\begin{itemize}
    \item Section \texttt{[settings]} complete et coherente avec la configuration de partie.
\end{itemize}
\end{besoinbox}

\subsection*{Sous besoins 3. Serialiser la section \texttt{[game]}}
\begin{besoinbox}
\textbf{Description :}
Ecrire l'etat du plateau, les racks et les scores de chaque joueur.

\textbf{Fonctions internes :}
\begin{verbatim}
saveBoard(PrintWriter writer, Board board)
serializeRack(Player p)
\end{verbatim}

\textbf{Format produit :}
\begin{itemize}
    \item Premiere ligne : index du joueur courant (1-based).
    \item Plateau ligne par ligne : lettre majuscule ou tiret \texttt{-} pour une case vide.
    \item \texttt{rack-N: LETTRES} — reserve de lettres du joueur N.
    \item \texttt{score-N: valeur} — score courant du joueur N.
    \item \texttt{language en} — langue de la partie (non dans \texttt{[settings]}).
\end{itemize}

\textbf{Test unitaire :}
\begin{itemize}
    \item \textbf{Cas 1 :} Entree : plateau avec mot \texttt{COUNT} en \texttt{g5h}. Attendu : la ligne \texttt{g} contient \texttt{C O U N T} aux bonnes colonnes.
    \item \textbf{Cas 2 :} Entree : plateau entierement vide. Attendu : toutes les cases sont representees par \texttt{-}.
    \item \textbf{Cas 3 :} Entree : joueur 1 score 15 / joueur 2 score 10. Attendu : presence de \texttt{score-1: 15} et \texttt{score-2: 10}.
\end{itemize}
\end{besoinbox}

\subsection*{Sous besoins 4. Serialiser la section \texttt{[history]}}
\begin{besoinbox}
\textbf{Description :}
Ecrire les coups joues dans l'ordre chronologique au format defini par F25.

\textbf{Fonctions internes :}
\begin{verbatim}
saveHistory(PrintWriter writer, Game game, List<Move> history)
readFullWordFromBoard(Board board, Move move)
convertPointToCoord(Point p)
\end{verbatim}

\textbf{Formats de coup supportes :}
\begin{itemize}
    \item \texttt{N pass} — le joueur passe son tour.
    \item \texttt{N exchange ABC} — le joueur echange des lettres.
    \item \texttt{N g5h WORD} ou \texttt{N e6v WORD} — coup normal horizontal ou vertical.
\end{itemize}

\textbf{Test unitaire :}
\begin{itemize}
    \item \textbf{Cas 1 :} Entree : joueur 1 joue \texttt{COUNT} en \texttt{g5h}, joueur 2 joue \texttt{PHOTO} en \texttt{e6v}. Attendu : presence de \texttt{1 g5h COUNT} et \texttt{2 e6v PHoTO} en tete d'historique.
    \item \textbf{Cas 2 :} Entree : aucun coup joue (partie neuve). Attendu : section \texttt{[history]} presente, sans ligne de coup.
    \item \textbf{Cas 3 :} Entree : coup de type \texttt{pass}. Attendu : format \texttt{N pass} ecrit correctement.
    \item \textbf{Cas 4 :} Entree : coup de type \texttt{exchange}. Attendu : format \texttt{N exchange LETTRES} ecrit correctement.
\end{itemize}
\end{besoinbox}

\subsection*{Sous besoins 5. Charger une partie depuis un fichier}
\begin{besoinbox}
\textbf{Description :}
Parser le fichier de sauvegarde et reconstruire un objet \texttt{Game} complet (joueurs, plateau, racks, scores, tour courant et historique).

\textbf{Fonction associee :}
\begin{verbatim}
Game loadGame(String filePath) throws Exception
\end{verbatim}

\textbf{Actions :}
\begin{itemize}
    \item Prelecture du fichier pour detecter le mode super-scrabble et creer le bon plateau (15$\times$15 ou 21$\times$21).
    \item Lecture sequentielle et parsing des sections \texttt{[settings]}, \texttt{[game]} et \texttt{[history]}.
    \item Reconstruction des joueurs humains et IA avec leurs parametres.
    \item Reconstruction de l'historique complet des coups.
\end{itemize}

\textbf{Gestion d'erreurs :}
\begin{itemize}
    \item Toute erreur de parsing est remontee sous la forme \texttt{Format error at line X: <detail>}.
    \item La CLI et la GUI affichent ensuite un message utilisateur lisible.
    \item L'etat courant du jeu n'est jamais ecrase en cas d'echec.
\end{itemize}

\textbf{Test unitaire :}
\begin{itemize}
    \item \textbf{Cas 1 :} Entree : fichier valide et complet. Attendu : \texttt{Game} non nul avec plateau, scores et historique coherents (\texttt{GameLoaderTest.java}).
    \item \textbf{Cas 2 :} Entree : plateau 21$\times$21 (super-scrabble). Attendu : detection correcte et creation du bon plateau.
    \item \textbf{Cas 3 :} Entree : fichier contenant \texttt{score-1: abc}. Attendu : exception \texttt{Format error at line X: ...}.
    \item \textbf{Cas 4 :} Entree : section \texttt{[game]} manquante. Attendu : exception avec message explicite.
    \item \textbf{Cas 5 :} Entree : coup \texttt{exchange} dans l'historique. Attendu : coup restaure correctement.
    \item \textbf{Cas 6 :} Entree : fichier inexistant. Attendu : exception fichier introuvable.
\end{itemize}
\end{besoinbox}

\subsection*{Sous besoins 6. Parser et ignorer les commentaires}
\begin{besoinbox}
\textbf{Description :}
Supprimer les commentaires avant le parsing tout en conservant une numerotation de ligne fiable pour les messages d'erreur, conformement a F22.

\textbf{Fonction interne :}
\begin{verbatim}
processLineComments(String line)
\end{verbatim}

\textbf{Comportement :}
\begin{itemize}
    \item Supprime tout ce qui suit un \texttt{\#} sur une ligne (commentaire en ligne).
    \item Ignore les blocs \texttt{\{ ... \}} potentiellement multi-lignes via un etat \texttt{isInBlockComment}.
    \item Continue le parsing ligne par ligne en maintenant le compteur de lignes.
\end{itemize}

\textbf{Test unitaire :}
\begin{itemize}
    \item \textbf{Cas 1 :} Entree : \texttt{debug=true \# activer le debug}. Attendu : \texttt{debug=true} apres suppression du suffixe commente.
    \item \textbf{Cas 2 :} Entree : bloc \texttt{\{ commentaire\textbackslash nsur 2 lignes \}} puis \texttt{language=fr}. Attendu : bloc supprime, \texttt{language=fr} conserve.
    \item \textbf{Cas 3 :} Entree : fichier sans commentaire. Attendu : contenu strictement identique en sortie.
\end{itemize}
\end{besoinbox}

% ------------------------------------------------------------------------------
\subsection*{Intégration dans l'architecture du projet}

\begin{table}[h!]
    \centering
    \renewcommand{\arraystretch}{1.3}
    \begin{tabular}{@{}lp{3.5cm}p{7cm}@{}}
        \toprule
        \textbf{Module} & \textbf{Responsabilité} & \textbf{Détail} \\
        \midrule
        \texttt{savefiles} & \texttt{SaveManager.java} \newline \texttt{GameLoader.java}
          & \texttt{SaveManager} : sérialisation de la partie. \newline \texttt{GameLoader} : parsing et reconstruction du \texttt{Game}. \\
        CLI & \texttt{GameControllerAux.java}
          & \texttt{save FILE} $\rightarrow$ \texttt{SaveManager.saveGame} \newline \texttt{load FILE} $\rightarrow$ \texttt{GameLoader.loadGame} \newline Erreurs affichées via \texttt{cliView.displayError} \\
        GUI & \texttt{ScrabbleGui.java}
          & Menu \textit{Save} $\rightarrow$ \texttt{SaveManager.saveGame} \newline Menu \textit{Load} $\rightarrow$ \texttt{GameLoader.loadGame} \newline Rebind du contrôleur et rafraîchissement UI après chargement. \\
        Configuration F2 & \texttt{ConfigLoader.java}
          & Lit \texttt{.scrabblerc} dans le HOME. \newline Note : le chemin par défaut de save/load depuis \texttt{.scrabblerc} n'est pas encore pris en charge si \texttt{FILE} est absent. \\
        \bottomrule
    \end{tabular}
\end{table}

% ------------------------------------------------------------------------------
\subsection*{Scénario utilisateur}

\begin{enumerate}
    \item L'utilisateur joue plusieurs coups en CLI ou via l'interface graphique.
    \item En CLI, il tape : \texttt{save ma\_partie.scrabble}
    \item \texttt{SaveManager} écrit \texttt{[settings]}, \texttt{[game]} puis \texttt{[history]} dans le fichier.
    \item Le shell affiche un message de succès à l'utilisateur.
    \item L'utilisateur relance le programme et tape : \texttt{load ma\_partie.scrabble}
    \item \texttt{GameLoader} reconstruit le \texttt{Game} : joueurs, plateau, tour courant et historique.
    \item Si une ligne est invalide, le programme affiche \texttt{Format error at line X: ...} et n'écrase pas l'état courant.
\end{enumerate}

%------------------------------------------------------------------------------
% BACK MATTER
%------------------------------------------------------------------------------
\appendix

\section{Section en annexe}

Lorem ipsum dolor sit amet, consectetur adipiscing elit. Sed do eiusmod tempor
incididunt ut labore et dolore magna aliqua. Ut enim ad minim veniam, quis
nostrud exercitation ullamco laboris nisi ut aliquip ex ea commodo consequat.
Duis aute irure dolor in reprehenderit in voluptate velit esse cillum dolore eu
fugiat nulla pariatur. Excepteur sint occaecat cupidatat non proident, sunt in
culpa qui officia deserunt mollit anim id est laborum.

\end{document}